\subsubsection{Enriched readers and fiber memory}
\label{act21:sec:ch-memoria-fibra}

The difference between history and value admits an operational criterion. At a
finite level \(H_k\), let \(q_k:H_k\to Z_k\) be a reader and let
\(M_k:H_k\to\mathcal M_k\) be the preserved record. The enriched reader
is \(\widehat q_k=(q_k,M_k)\). Memory distinguishes states within a single
fiber of the reader; its sufficiency is verified on the domain under consideration.
This formulation incorporates into the article itself the complete development of the
fiber information used below.

\begin{proposition}[Recoverability at a finite level]
\label{act21:prop:ch-recuperabilidad-fibra}
An inverse of \(\widehat q_k\) exists on its image if and only if
\(\widehat q_k\) is injective. For a random variable \(Z\) taking
values in \(H_k\) with distribution \(\mu_k\),
\[
 \mathsf H(Z)=\mathsf H(q_k(Z))+\mathsf H(Z\mid q_k(Z)).
\]
If \(\widehat q_k\) is injective on the support of \(Z\), then
\[
 \mathsf H(Z\mid q_k(Z),M_k(Z))=0.
\]
Here \(\mathsf H(Z)=-\sum_z\Pr(Z=z)\log\Pr(Z=z)\), with
\(0\log0=0\), is finite entropy in nats; conditional entropy is
the average entropy of the conditional distributions.
\end{proposition}

\begin{proof}
An inverse recovers a unique antecedent, which requires injectivity;
conversely, injectivity makes it possible to assign to each expressed pair its unique
antecedent. Since \(q_k(Z)\) is a deterministic function of \(Z\),
grouping the sum defining \(\mathsf H(Z)\) by fibers of \(q_k\) yields the
chain rule. If the pair \((q_k(Z),M_k(Z))\) distinguishes the support,
each corresponding conditional distribution is concentrated on a single
state, and its entropy is zero.
\end{proof}

Not publishing a coordinate and erasing the record that permits its recovery are
distinct operations. This result does not identify entropy with the cardinality
of the continuum or introduce a thermal interpretation: its domain, reader, and
distribution are precisely those just specified.

\subsubsection{Equivariance of the measure under tree symmetries}
\label{act21:sec:ch-naturalidad-medida}

\begin{proposition}[Joint transport of the tree and its measure]
\label{act21:prop:ch-naturalidad-medida}
Let \(g=(g_k)_{k\geq0}\) be a collection of bijections \(g_k:H_k\to H_k\)
that commutes with truncations, fixes the root, and bijectively transports
admissible emissions while preserving their multiplicities. Then
\[
 d(g_k\zeta)=d(\zeta),\qquad
 (g_k)_*\mu_k=\mu_k,\qquad (g_\infty)_*\mu=\mu.
\]
More generally, if \(\mu_\zeta\) is the law of continuations of a
history \(\zeta\in H_k\), then
\[
 (g_\infty)_*\mu_\zeta=\mu_{g_k\zeta}.
\]
\end{proposition}

\begin{proof}
The bijection of emissions preserves their number with multiplicities. The
uniform seed distribution is likewise invariant. Each path
and its image therefore receive the same product of local weights
\(1/d(\zeta)\). Equality of measures follows at every level and on
every cylinder. Compatibility with truncations defines \(g_\infty\),
and uniqueness of the measure determined by cylinders proves equality
in the limit. The same argument, beginning at \(\zeta\), proves the
assertion about continuations.
\end{proof}

This establishes the naturality of the tree and its measure: a coherent change
of representation transports the weights and does not require choosing them anew.
Preservation of emissions and multiplicities is an essential hypothesis of
the proposition, not a datum added after the construction.
